%% Section 2 — Methods
%% CMPB submission: TRUST-CT
%% Locked values drawn from phase3_lock.json, phase4_lock.json,
%% phase5_v3_lock.json, phase6_c0_lock.json, ablation_summary.json.
%% DO NOT edit numerical values without updating numerical_provenance.csv.

\section{Methods}
\label{sec:methods}

\subsection{Study design and TRUST-CT workflow}
\label{sec:design}

This was a retrospective, multi-cohort computational evaluation of
a closed-loop federated learning (FL) framework for budget-constrained
medication-adherence incentive allocation with signed governance logging.
The pipeline follows the sequence: provider-local clinical data
$\to$ local model training $\to$ engagement-policy evaluation
$\to$ trial-calibrated availability replay $\to$ federated aggregation
$\to$ audit logging (Figure~\ref{fig:workflow}).
Five roles are defined: (i)~\emph{providers} retain raw records locally;
(ii)~a~\emph{coordinator} aggregates parameter updates;
(iii)~a~\emph{policy module} ranks participants for incentive allocation;
(iv)~an~\emph{attacker} may submit adversarial updates;
(v)~an~\emph{audit verifier} checks log integrity.
Raw clinical records were not transmitted to the coordinator.
Secure aggregation and distributed consensus were not implemented;
the coordinator observes individual client updates in plaintext.

\subsection{Clinical cohorts and outcomes}
\label{sec:cohorts}

Three cohorts provided complementary experimental roles
(Table~\ref{tab:cohorts}).

\paragraph{Cimas (principal prediction and closed-loop cohort).}
A retrospective hypertension medication-adherence dataset comprising
claims from 12{,}649 participants ($K=23$ primary providers,
each with $\geq 100$ participants) was extracted under a
six-month landmark design (January--June 2022 observation window;
refill-continuity outcome assessed July--December 2022).
The binary outcome $Y_\text{adh}=1$ denoted adherence
($\geq 5$ refill months; prevalence 42.4\%).
For the long-tail evaluation, the risk-positive orientation
$Y_\text{risk}=1-Y_\text{adh}$ was used, with a long-tail
prevalence of 0.677.
Date fields followed D/M/YYYY format (confirmed from birth-date records).

\paragraph{BETTER-BP (policy evaluation and replay calibration cohort).}
The BETTER-BP pragmatic randomized controlled trial enrolled
participants at three safety-net clinics under a 2:1
(intervention:control) allocation.
Of 402 randomized participants, two FL providers (Sites A, $n=312$;
Site B, $n=90$; Site C excluded, $n=0$) contributed to policy
evaluation and closed-loop calibration.
% W02a fix: outcome is visit_6m_attended (Y=1=attended), confirmed from
%   cohort_flow.json and run_phase5_v3.py line 230. Corrected from
%   "12-month non-attendance" to "6-month visit attendance".
The primary outcome was six-month visit attendance ($Y=1$=attended).
% AUTHOR_PENDING: confirm precise 6-month attendance count and ethics
%   approval reference. Also confirm ClinicalTrials.gov NCT04114669
%   applies to this cohort and consent procedure for secondary analysis.

\paragraph{eICU (adversarial security and privacy benchmark).}
The eICU Collaborative Research Database (demo version)
provided a multi-centre critical-care benchmark~\cite{TODO:pollard2018eicu}.
The cohort comprised 2{,}520 hospital stays across
$K=8$ hospital groups; the outcome was in-hospital mortality
(126 positives; 5.0\% prevalence).
This cohort was used exclusively for adversarial robustness and
privacy-leakage evaluation; it was not used for clinical effectiveness claims.

\begin{table}[ht]
\centering
\caption{Clinical cohorts, outcomes and experimental roles.}
\label{tab:cohorts}
\begin{tabular}{@{}llrlll@{}}
\hline
Cohort & $N$ & $K$ & Primary outcome & Role & Evaluation \\
\hline
Cimas     & 12{,}649 & 23 & Refill adherence ($\geq$5 months) & FL prediction; closed loop & LOCO / landmark \\
BETTER-BP & 402      & 2  & 6-month visit attendance ($Y=1$=attended)  & Policy est.; replay calib. & Trial secondary \\
eICU      & 2{,}520  & 8  & In-hospital mortality (5.0\%)     & Security / privacy bench.  & Train/test split \\
\hline
\end{tabular}
\end{table}

\subsection{Preprocessing, splitting and feature provenance}
\label{sec:preproc}

For the eICU security benchmark, preprocessing followed a
strictly leakage-safe order:
(1)~construct the binary outcome and eligible cohort;
(2)~stratified train/test split ($20\%$ held out; random state 7);
(3)~fit a median imputer on training rows only;
(4)~transform training and test rows using the frozen imputer;
(5)~compute mutual information (MI) using training features and
training labels only;
(6)~select the top-60 features by MI rank;
(7)~freeze and hash the selected feature list
    (SHA-256: \texttt{4f4a66036e869328b354fb528e8dface\ldots}).
No test labels, test rows, or full-cohort imputation contributed
to feature selection.
The implementation script was independently hashed
(SHA-256: \texttt{93200b7d30039a9d\ldots}).

\subsection{Centralized and federated prediction}
\label{sec:fl}

Each provider $k$ minimized the local empirical risk:
\begin{equation}
  F_k(\mathbf{w}) = \frac{1}{n_k}\sum_{i=1}^{n_k}
  \ell\!\left(\mathbf{w};\, \mathbf{x}_{ki},\, y_{ki}\right),
  \label{eq:local_obj}
\end{equation}
where $\ell$ is logistic loss and $n_k$ is the local sample size.
The global objective is $F(\mathbf{w}) = \sum_k \tfrac{n_k}{N}F_k(\mathbf{w})$.

Local updates used mini-batch stochastic gradient descent
($N_\text{local}=5$ epochs, learning rate $\eta=0.02$,
balanced class weights).
Five aggregation strategies were evaluated:
(i)~\emph{FedAvg} (weighted average~\cite{TODO:mcmahan2017fedavg});
(ii)~\emph{FedProx} (proximal regularisation $\mu$);
(iii)~\emph{FedAvg-equal} (uniform weights);
(iv)~\emph{SCAFFOLD} (control variates~\cite{TODO:karimireddy2020scaffold});
(v)~\emph{FedAdam} (adaptive server update).
The Cimas experiment ran for 30 rounds with 23 provider-clients.
Leave-one-client-out (LOCO) cross-validation was used for evaluation.

The noninferiority estimand is $d = \text{AUROC}_\text{FL} -
\text{AUROC}_\text{central}$ with prespecified margin $\delta_0 = -0.02$.
FL was declared non-inferior when the one-sided 95\% lower
confidence bound (LCB) of $d$ exceeded $\delta_0$.
% W05a fix: N_BOOT=500 confirmed from run_phase5_v3.py line 50 (N_BOOT = 500).
CIs were obtained by provider-cluster and participant bootstrap
(500 resamples).

\subsection{Engagement-policy estimation}
\label{sec:policy}

% W02a fix: outcome is six-month visit attendance, not twelve-month.
The treatment was binary incentive eligibility; the outcome was
six-month visit attendance ($Y=1$=attended).
Doubly robust (DR) policy value estimation used the estimator:
\begin{equation}
  \hat{V}(\pi) = \frac{1}{n}\sum_{i=1}^{n}\!\left[
    \hat{\mu}_{\pi(X_i)}(X_i)
    + \frac{\mathbf{1}{\{A_i = \pi(X_i)\}}}{\hat{e}_{A_i}(X_i)}
    \bigl(Y_i - \hat{\mu}_{A_i}(X_i)\bigr)
  \right],
  \label{eq:dr}
\end{equation}
where $\hat{\mu}_a(\cdot)$ is a cross-fitted outcome regression and
$\hat{e}_a(\cdot)$ is the propensity model.
The BETTER-BP randomization allocated participants 2:1
(intervention:control; known propensity $e=0.667$).
Five-fold, five-seed out-of-fold (OOF) cross-fitting was used.
Four policies were compared at budget $B=\tfrac{2}{3}$:
$\pi_0$ (no incentive), $\pi_\text{rand}$ (random),
$\pi_\text{risk}$ (FL nonattendance-risk ranking), and
$\pi_\text{up}$ (cross-fitted DR uplift ranking).
Causal assumptions of consistency, conditional exchangeability
(by trial randomization), and positivity were assessed.
No prospective policy validation was performed.

\subsection{Trial-calibrated closed-loop replay}
\label{sec:replay}

Four policy scenarios (C0--C3) were evaluated in a trial-calibrated
simulation using attendance and allocation random streams generated
independently from common random numbers (CRN) shared across policies.

Simulated attendance probabilities were recalibrated on the logit scale:
\begin{equation}
  p_{ia}^\text{cal} = \sigma\!\left[
    \sigma^{-1}\!\left(\hat{\mu}_a(X_i)\right) + \delta_a
  \right],
  \label{eq:calib}
\end{equation}
where $\sigma$ is the sigmoid function and $\delta_a$ is selected by
bisection (Brent's method) so that
$\mathbb{E}[p_{i0}^\text{cal}]=0.815$ and
$\mathbb{E}[p_{i1}^\text{cal}]=0.883$,
reproducing $\hat{V}(C1, B=\tfrac{2}{3})=0.8603$ from Phase~4.
% W04 Option A reframing: the FL model is refitted independently in each
% outer round from the round-specific attending subset. The outer rounds
% do NOT carry model state from round to round; each refitting starts from
% the warm-initialised weights and runs N_FL=5 inner aggregation rounds.
% Language about "sequential" or "continuous" FL learning across outer
% rounds has been removed; this is policy-stratified repeated refitting.
The Cimas cold-start experiment ran for $T=30$ outer replay rounds with a
fixed 20\% stratified holdout, $5$~model seeds, and $10$~CRN draws per seed.
At each outer round, FL was refitted on the attending participant subset
($N_\text{local}=5$ inner aggregation rounds); performance was measured on
the fixed holdout.
Provider volume was matched within-provider across scenarios.
Performance was summarised by normalized area under the learning
curve (nAULC) over the 30 replay rounds.
The operational materiality threshold was $|\Delta\,\text{nAULC}| < 0.001$.
This replay is a trial-calibrated simulation; it does not constitute
independent causal or prospective clinical validation.

\subsection{Threat model and adversarial federated learning}
\label{sec:threat}

Seven adversarial roles were defined (Supplementary~S2):
malicious clinical site; colluding sites; honest-but-curious coordinator;
external attacker; and three governance roles.
The coordinator observes client updates in plaintext; secure aggregation
was not implemented.
Seven attack families were evaluated:
(i)~\emph{sign flipping};
(ii)~\emph{label flipping};
(iii)~\emph{scaled poisoning};
(iv)~\emph{model replacement} (scale $K/f$);
(v)~\emph{A Little Is Enough} (ALIE; colluding gradient statistics~\cite{TODO:baruch2019alie});
(vi)~\emph{random Gaussian perturbation};
(vii)~\emph{targeted tabular backdoor} (trigger feature 0, value 3.0, target label 0).
Experiments varied $f$ malicious clients: eICU $f \in \{1,2,3\}$
(12.5\%, 25.0\%, 37.5\% of $K=8$);
Cimas $f \in \{2,5,7\}$ (8.7\%, 21.7\%, 30.4\% of $K=23$).
All experiments used 15 federated rounds and 5 random seeds.

\subsection{Robust aggregation and feasibility}
\label{sec:robust}

Five aggregation rules were evaluated.
\emph{Coordinate-wise median}:
$[\text{agg}(\{\boldsymbol{\delta}_k\})]_j = \mathrm{median}_k\, \delta_{kj}$.
\emph{Trimmed mean}: trims the $\lfloor f \rfloor$ smallest and
largest values per coordinate before averaging.
\emph{Krum}: selects the update minimizing the sum of distances to its
$K-f-2$ nearest neighbours.
Krum requires $K > 2f+2$; for eICU with $K=8$, $f=3$,
this condition fails ($8 \not> 8$).
The 35 Krum-inapplicable configurations (eICU, $f=3$, adversarial attacks)
used FedAvg as a diagnostic fallback
(\texttt{aggregator\_effective = fedavg},
\texttt{run\_status = not\_applicable\_diagnostic\_fallback})
and are excluded from Krum inference.

\subsection{Empirical leakage evaluation}
\label{sec:leakage}

Each client update was perturbed before transmission:
\begin{equation}
  \tilde{\boldsymbol{\Delta}}_k =
  \mathrm{clip}_C(\boldsymbol{\Delta}_k) +
  \mathcal{N}\!\left(\mathbf{0},\, (\alpha C)^2 \mathbf{I}\right),
  \label{eq:pert}
\end{equation}
where $C=1.0$ is the clipping norm and $\alpha \in \{0,0.01,0.05,0.10,0.20\}$
is an empirical noise multiplier.
No $(\varepsilon,\delta)$-differential privacy accounting was performed;
the mechanism is described as
\emph{utility-oriented clipped Gaussian update perturbation providing
empirical leakage mitigation}.

Three reconstruction settings were differentiated:
(1)~\emph{analytic single-step canary}: batch size 1, one gradient step,
logistic regression; the exact recover $\hat{\mathbf{x}} = \nabla_\mathbf{w}\ell
/ \nabla_b\ell$ is valid and clipping-invariant;
(2)~\emph{optimistic fresh-gradient reconstruction} (\texttt{recon\_single}):
the coordinator observes the gradient at the global model for batches of 1,
4, and 16 samples --- an optimistic attacker scenario;
(3)~\emph{multi-step delta approximation}: the same formula applied to the
coordinator-observed multi-epoch FedAvg update, which is invalid
(Supplementary~S3).
Leakage metrics included cosine similarity, normalised RMSE,
feature-sign recovery rate, and membership-inference (MI) attack AUROC.
MI attack advantage is $\text{AUROC}-0.5$; effect size is Cohen's
$h = 2\arcsin\!\sqrt{p_1} - 2\arcsin\!\sqrt{0.5}$.

% W11a/b/c fix: oracle.py uses HMAC-SHA256 (symmetric key), NOT public-key.
% Removed: signature/non-repudiation/public-key language.
% Removed: "Signed" from subsection title.
% Added: explicit statement that HMAC does not provide non-repudiation.
% Added: 64-bit truncation disclosure for chain links.
\subsection{Hash-linked audit log with HMAC authorization}
\label{sec:auditlog}

Each governance event was serialised to a canonical record
containing: round index, provider identifier, model hash, attendance
assertion, action taken, and timestamp.
Two cryptographic fields were computed per record $t$:
\begin{align}
  d_t &= H(\text{canonical}(e_t)), \label{eq:digest} \\
  s_t &= \text{HMAC-SHA256}(\mathit{key},\, d_t), \label{eq:auth} \\
  c_t &= H_{\!64}(c_{t-1} \,\|\, d_t \,\|\, s_t), \label{eq:chain}
\end{align}
where $H$ is SHA-256, $H_{\!64}$ denotes the first 64 bits of the
SHA-256 digest (used for chain efficiency), and $\mathit{key}$ is a
shared authorization key per provider.
This mechanism provides \emph{HMAC authorization} (source authentication)
and \emph{sequential hash-chain integrity} (tamper and reordering detection);
it does not provide non-repudiation, which requires an asymmetric signing key
inaccessible to the verifier.
Verification checked HMAC consistency, event-digest validity,
and sequential previous-hash linkage, detecting tampering,
reordering, replay, deletion, substitution, and stale attestations.
Seven fault classes were injected across 10 round positions
and 5 seeds (350 injections total).
The log authenticates the recorded source and detects subsequent
alteration, but cannot establish the truth of an incorrect event
submitted by an authorized source.

\subsection{Statistical analysis and reproducibility}
\label{sec:stats}

Primary metrics were AUROC, PR-AUC (risk-positive orientation for Cimas),
balanced accuracy, MCC, and DR policy value.
% W10a fix: BSR renamed to triggered_target_class_rate_all; denominator clarified.
% W05a fix: bootstrap resamples corrected to 500 (confirmed from N_BOOT in code).
Security metrics included attack degradation ($\Delta$AUROC paired over seeds),
defense recovery, triggered target-class rate
(proportion of all triggered test rows predicted as the target class,
\emph{not} restricted to non-target-class rows; see Section~\ref{sec:threat}),
reconstruction cosine similarity, feature-sign recovery, and MI AUROC.
Paired t-tests with 95\% CIs (two-sided, $\alpha=0.05$) were used for
attack-defense comparisons ($n=5$ seeds).
Bootstrap CIs (500 resamples, percentile method) were computed for
AUROC, PR-AUC, and policy-value contrasts.
MI attack CIs used 500-resample bootstrap at the per-seed level.
The noninferiority margin was $\delta_0=-0.02$ AUROC;
the utility materiality threshold was $|\Delta\text{AUROC}|<0.02$;
the closed-loop operational materiality threshold was
$|\Delta\,\text{nAULC}|<0.001$.
All experiments used fixed random seeds ([7,11,19,23,37]);
artifact hashes are reported in Supplementary~S1.

\subsection{Provider-level performance disparity analysis (E6)}
\label{sec:e6_design}

A four-condition warm/cold de-confound experiment was conducted on the Cimas
primary cohort to isolate federation benefit from warm weight initialisation.
The four conditions (\emph{cold-FedAvg}, \emph{cold-local}, \emph{warm-local},
\emph{warm-FedAvg}) shared identical test sets (within-provider stratified 80/20
splits fixed by seed), preprocessing (per-provider median imputer and
StandardScaler fitted on training rows only), feature set (22~numeric columns),
local epochs ($N_\text{local}=5$), learning rate ($\eta=0.02$), and
class weights (balanced).
Warm conditions were initialised from the cold-FedAvg round-30 global weights;
warm-FedAvg continued aggregation from that point and served as an exact
reproducibility gate (not an independent arm).
Per-provider AUROC was evaluated at each round using each provider's
locally-fitted preprocessor applied to its held-out test rows.
Inference used seed-level paired $t$-tests ($n=5$ seeds, df=4) with
$5{,}000$-resample seed-cluster bootstrap 95\%~CIs.
Provider-level disparity was quantified as IQR and worst-provider AUROC
across $K=23$~providers; these statistics describe provider-level performance
equity and do not constitute demographic or individual fairness assessments.
Full design and output files are described in Supplementary~S5.
